% from ega2-errata-addenda-fr.tex
\oldpage[II]{219}
\markboth{ERRATA ET ADDENDA, LISTE 1}{ERRATA ET ADDENDA, LISTE 1}
de \emph{morphisme d'espaces annelés
en anneaux locaux}, c'est toujours d'un tel morphisme qu'il s'agira;
avec cette définition des morphismes, il est clair que les espaces
annelés en anneaux locaux forment une catégorie; pour deux objets
$X,Y$ de cette catégorie, $\Hom(X,Y)$ désignera donc l'ensemble des
morphismes d'espaces annelés en anneaux locaux de $X$ dans $Y$
(ensemble noté $\mathfrak L$ dans (1.8.1)); lorsque nous aurons à
considérer l'ensemble des morphismes d'espaces annelés de $X$ dans
$Y$, nous le noterons $\Hom_{\mathrm{an}}(X,Y)$ pour éviter toute
confusion. L'application (1.8.1.1) s'écrit donc
\begin{equation}
\label{ega2:addendum:I.1.8.2.1-fr}
  \rho:\Hom_{\mathrm{an}}(X,Y)
  \to\Hom(\Gamma(Y,\sh{O}_Y),\Gamma(X,\sh{O}_X)),
\tag{1.8.2.1}
\end{equation}
et sa restriction
\begin{equation}
\label{ega2:addendum:I.1.8.2.2-fr}
  \rho':\Hom(X,Y)
  \to\Hom(\Gamma(Y,\sh{O}_Y),\Gamma(X,\sh{O}_X))
\tag{1.8.2.2}
\end{equation}
est une application fonctorielle en $X$ et $Y$ dans la catégorie des
espaces annelés en anneaux locaux.
\end{env}

\begin{corollary}[1.8.3]
\label{ega2:addendum:I.1.8.3-fr}
Soit $(Y,\sh{O}_Y)$ un espace annelé en anneaux locaux. Pour que $Y$
soit un schéma affine, il faut et il suffit que pour tout espace annelé
en anneaux locaux $(X,\sh{O}_X)$, l'application (1.8.2.2) soit
bijective.
\end{corollary}

La proposition (1.8.1) montre que la condition est nécessaire.
Inversement, si on la suppose vérifiée et si l'on pose
$A=\Gamma(Y,\sh{O}_Y)$, il résulte de l'hypothèse et de (1.8.1) que les
foncteurs $X\mapsto\Hom(X,Y)$ et
$X\mapsto\Hom(X,\Spec(A))$, de la catégorie des espaces annelés en
anneaux locaux dans celle des ensembles, sont isomorphes. On sait que
cela entraîne l'existence d'un isomorphisme canonique
$X\to\Spec(A)$ (cf.~0, 8).

\begin{env}[1.8.4]
\label{ega2:addendum:I.1.8.4-fr}
Soit $S=\Spec(A)$ un schéma affine; désignons par $(S',A')$ l'espace
annelé dont l'espace sous-jacent est réduit à un point et le faisceau
structural $A'$ est le faisceau (nécessairement simple) sur $S'$ défini
par l'anneau $A$. Soit $\pi:S\to S'$ l'unique application de $S$ dans
$S'$; notons d'autre part que pour tout ouvert $U$ de $S$, on a une
application canonique
$\Gamma(S',A')=\Gamma(S,\sh{O}_S)\to\Gamma(U,\sh{O}_S)$ qui définit
donc un $\pi$-morphisme $\iota:A'\to\sh{O}_S$ de faisceaux d'anneaux.
On a ainsi défini canoniquement un morphisme d'espaces annelés
$i=(\pi,\iota):(S,\sh{O}_S)\to(S',A')$. Pour tout $A$-module $M$, nous
désignerons par $M'$ le faisceau simple sur $S'$ défini par $M$, qui
est évidemment un $A'$-Module. Il est clair que l'on a
$i_*(\widetilde M)=M'$ (I, 1.3.7).
\end{env}

\begin{lemma}[1.8.5]
\label{ega2:addendum:I.1.8.5-fr}
Avec les notations de (1.8.4), pour tout $A$-module $M$, le
$\sh{O}_S$-homomorphisme canonique fonctoriel (0, 4.4.3.3)
\begin{equation}
\label{ega2:addendum:I.1.8.5.1-fr}
  i^*(i_*(\widetilde M))\to\widetilde M
\tag{1.8.5.1}
\end{equation}
est un isomorphisme.
\end{lemma}

En effet, les deux membres de (1.8.5.1) sont exacts à droite (le
foncteur $M\mapsto i_*(\widetilde M)$ étant évidemment exact) et
commutent aux sommes directes; en considérant $M$ comme conoyau d'un
homomorphisme $A^{(I)}\to A^{(J)}$, on se ramène à prouver le lemme
dans le cas $M=A$, et il est évident dans ce cas.

\begin{corollary}[1.8.6]
\label{ega2:addendum:I.1.8.6-fr}
Soient $(X,\sh{O}_X)$ un espace annelé, $u:X\to S$ un morphisme
d'espaces
